\documentclass[10pt,a4paper]{amsart}
\usepackage[margin=19mm]{geometry}
\usepackage{amsmath,amssymb,mathtools}
\usepackage[scr=rsfs,cal=euler]{mathalfa}
\usepackage{xfrac}
\usepackage[unicode,hidelinks,bookmarks=false]{hyperref}
\newcommand{\ie}{i.e.\ }
\newcommand{\cf}{cf.\ }
\newcommand{\resp}{respectively\ }
\newcommand{\Subsection}{\subsection}
\providecommand{\nobreakdash}{\nobreak\hbox{-}}
\setlength{\emergencystretch}{2em}
\multlinegap0pt
\usepackage{amssymb}
\usepackage{subfig}
\usepackage{tikz}
\usepackage[shortlabels]{enumitem}
\usepackage{xcolor}
\usepackage{bm}
\newtheorem{lemma}{Lemma}
\newtheorem{theorem}{Theorem}
\newtheorem{corollary}{Corollary}
\title{Orbit equivalence of Cantor minimal systems}
\author{Su Gao, Ruiwen Li, Yiming Sun}
\date{Source: arXiv:2601.00179v1 (CC BY 4.0)}
\begin{document}
\maketitle

\noindent\emph{Source and licence.} Orbit equivalence of Cantor minimal systems. Version arXiv:2601.00179v1 (CC BY 4.0), distributed by arXiv under the Creative Commons Attribution 4.0 International licence, http://creativecommons.org/licenses/by/4.0/, as recorded in the arXiv licence metadata for this exact version and shown on the abstract page https://arxiv.org/abs/2601.00179v1. Full licence notice: \texttt{licenses/arxiv-2601.00179-CC-BY-4.0.txt}.\par\smallskip

\noindent\textit{Editorial scope.} Counted unit: the article's theorem theorem (source0.tex lines 975--976). The article's complete original proof of it is delivered with it (source0.tex lines 978--996, one \texttt{proof} environment of 19 lines). No mathematics is written, repaired or re-derived: the statement and the proof are the article's own lines, in the article's own notation, and no retained line is substituted or re-flowed.  Retained uncounted as the source-local context the proof needs: lines 902--903; lines 948--949; lines 959--960.  Nothing was omitted from any retained span, so the package carries no omission record.  No retained line contains a citation, so the wrapper carries no bibliography and no bibliography entry.  No retained line contains a figure environment, an image-inclusion command or drawing code, so no image is delivered and the package declares no asset dependency.  Editorial formatting only, in addition: an AMS \texttt{amsart} preamble that restates the article's own declarations and macros used by the retained lines, namely the environments \texttt{lemma}, \texttt{theorem}, \texttt{corollary} on separate counters, together with the standard packages those lines need.  The retained environments print in this document as Lemma 1, Theorem 1, Corollary 1 in order of appearance, whereas the article prints the same items with the numbering of its own full document.  The complete native source of the article as distributed in the arXiv e-print for this exact version is bundled unchanged as \texttt{sources/191977/source0.tex}.
\begin{lemma}\label{lem:NKL} $\mathcal{R}_{N,K,L}=\varnothing$ for $L>N$.
\end{lemma}

\begin{theorem}\label{thm:rn} For any $N\geq 2$, $R_N$ is virtually countable.
\end{theorem}

\begin{corollary}\label{cor:L1} For any $N\geq 2$ and $1\leq K\leq N$, $R_{N,K,1}$ is virtually hyperfinite.
\end{corollary}

\begin{theorem} For any $N\geq 2$ and $1\leq K\leq N$, $R_{N, K, 2}$ is Borel reducible to an orbit equivalence relation induced by a Borel action of a countable amenable group. In particular, $R_2$ is Borel reducible to an orbit equivalence relation induced by a Borel action of a countable amenable group. Consequently, $R_2$ is virtually amenable. 
\end{theorem}

\begin{proof} The key point is to note that if $(X, \varphi)\in \mathcal{R}_{N,K, L}$, then we may always choose $1^K$ to be an element of a maximally $\mathbb{Q}$-linearly independent subset of $\Gamma_\varphi$. Thus if $1\leq K\leq 2$ and $L=2$, then we may assume $1^K$ and $\vec{a}^\varphi\in \mathbb{R}^K$ form a maximally $\mathbb{Q}$-linearly independent subset of $\Gamma_\varphi$. 

Define an equivalence relation $Q_K$ on $\mathbb{R}^K\times 2^{\mathbb{Q}^2}$ as follows. For $\vec{a}, \vec{b}\in \mathbb{R}^K$ and $p,q\in 2^{\mathbb{Q}^2}$, 
$$ \mbox{ $(\vec{a}, p)$ and $(\vec{b}, q)$ are $Q_K$-equivalent} $$
if and only if there are $\phi\in\mbox{\rm Sym}(K)$ and $r_0, s_0\in \mathbb{Q}$ with $r_0\neq 0$ such that
$$ r_0\phi(\vec{a})+s_01^K=\vec{b} $$
and for any $r, s\in \mathbb{Q}$,
$$ p(r, s)=q(rr_0, rs_0+s). $$

Note that $Q_K$ is the orbit equivalence relation induced by a continuous action of the group
$$ G=(\mbox{\rm GL}(1, \mathbb{Q})\ltimes \mathbb{Q})\times \mbox{\rm Sym}(K), $$
where $\mathbb{Q}$ is the additive group and the semi-direct product is defined by
$$ (r, s)(r',s')=(rr', rs'+s) $$
for $r,r'\in\mathbb{Q}\setminus\{0\}$ and $s, s'\in\mathbb{Q}$. Since $G$ is solvable of rank $2$, it is amenable. 

By a similar argument as in the proof of Theorem~\ref{thm:rn}, we have that $R_{N, K, 2}$ is Borel reducible to $Q_K$. 

The second part of the theorem follows from Lemma~\ref{lem:NKL} and Corollary~\ref{cor:L1}.
\end{proof}

\end{document}
